\section{Embasamento Teórico}
\label{sec-embasamento}

Nesta seção, é apresentado os conceitos fundamentais de Teoria dos Grafos e de heurísticas aplicadas à resolução de problemas de Otimização Combinatória, os quais foram consolidados ao longo desta Iniciação Científica.

A Teoria dos Grafos, estuda estruturas compostas por vértices e arestas, e possui ampla aplicação em domínios como redes de comunicação, transporte e biologia. A representação dessas estruturas pode ser feita por listas ou matrizes de adjacência, facilitando a compreensão de conceitos centrais, como adjacência, conectividade, caminhos, ciclos e vizinhança, todos abordados e revisados em \cite{grafos:2017}. Neste trabalho, destaca-se o estudo de árvores e árvores geradoras. Uma árvore é um grafo que conecta todos os seus vértices sem formar ciclos, garantindo a existência de um único caminho entre quaisquer dois vértices. Já uma árvore geradora de um grafo $G$ é um subgrafo acíclico que conecta todos os vértices de $G$ com exatamente $|V| - 1$ arestas, sendo que a remoção de qualquer aresta resulta na desconexão da árvore.

Outro conceito relevante, durante o estudo atual, é a densidade de um grafo, métrica que avalia o grau de conectividade de sua estrutura. A densidade corresponde à razão entre o número de arestas existentes e o número máximo de arestas possíveis, variando entre 0 e 1. Valores próximos de 1 indicam grafos densos (ou completos), enquanto valores próximos de 0 caracterizam grafos esparsos (poucas conexões). 

A centralidade de grafos também desempenha papel essencial neste estudo, pois permite avaliar a importância relativa dos vértices. Essa medida classifica os vértices segundo diferentes critérios, como: (i) centralidade de grau, que avalia a conectividade direta; (ii) centralidade de proximidade, que considera a distância média em relação aos demais vértices; e (iii) centralidade de intermediação, que mensura a frequência de um vértice em caminhos mínimos. Neste trabalho, a análise se concentra em duas medidas específicas: a centralidade de grau e o \textit{PageRank}.  O \textit{PageRank}, foi proposto por \cite{PageRank:1998}, e atribui importância a um vértice com base não apenas na quantidade, mas também na qualidade de suas conexões, conferindo maior relevância aos vértices ligados a outros igualmente importantes. 

Em estudos sobre Otimização Combinatória, a combinação de heurísticas fundamentais, como algoritmos construtivos e busca local, é essencial para encontrar soluções eficientes em problemas complexos. Algoritmos construtivos geram soluções de forma incremental sem revisar escolhas anteriores, enquanto a busca local refina uma solução inicial ao explorar sua vizinhança. A união desses princípios deu origem a meta-heurísticas, métodos que equilibram a exploração global e a intensificação local do espaço de busca. A meta-heurística \textit{Multi-Start} é um exemplo clássico, pois realiza múltiplas reinicializações de busca local a partir de soluções geradas aleatoriamente, evitando que o algoritmo se prenda em ótimos locais. O GRASP (\textit{Greedy Randomized Adaptive Search Procedure}) é uma abordagem mais sofisticada que refina a ideia do \textit{Multi-Start}, combinando uma fase construtiva que usa critérios gulosos e aleatoriedade adaptativa com uma busca local para refinar as soluções obtidas. Essas técnicas, amplamente discutidas em \cite{metaheuristicas2003}, constituíram ferramentas centrais ao longo deste trabalho.
